\chapter*{Заключение}
\addcontentsline{toc}{chapter}{Заключение}
\label{ch:Заключение}

В рамках настоящей работы проведено исследование применимости мультиязычных трансформерных архитектур к задаче трёхклассовой классификации тональности текста в условиях клиентской браузерной среды с учётом ограничений WebAssembly, WebGPU и аппаратных ресурсов конечных устройств.

В ходе теоретического исследования выявлен комплекс системных ограничений среды JavaScript, исключающих интерактивный нейросетевой вывод в главном потоке без применения WebAssembly и Web Workers; для преодоления вычислительных ограничений браузерной среды применяется многоуровневая стратегия оптимизации, сочетающая дистилляцию знаний, динамическое квантование INT8 и аппаратное ускорение посредством WASM SIMD128, при которой дистилляция устраняет структурную избыточность на уровне архитектуры, квантование снижает ресурсоёмкость на уровне представления параметров, а WASM SIMD128 обеспечивает векторизованное исполнение операций на уровне среды выполнения.

Разработанная трёхуровневая архитектура клиентского SPA (React~+~Transformers.js~+~ONNX Runtime Web) реализует полный цикл вывода в изолированном потоке Web Worker; дообученная модель ayayousef/distilbert-multilingual-sentiment-finetuned опубликована на HuggingFace Hub, а квантование сократило объём файла с 541{,}4 до 135{,}8~МБ (редукция 74{,}9\%). Также реализован конвейер CI/CD с автоматическим развёртыванием приложения в Yandex Serverless Containers.

Экспериментальная валидация на 22\,500 образцах в трёх языках подтвердила центральную гипотезу: конфигурация DistilmBERT FT INT8 достигает Overall F1-macro~=~0{,}7224 при задержке 94{,}91~мс/образец и индексе Browser Score~=~0{,}998, причём деградация качества при квантовании составила лишь $\Delta$F1~=~$-$0{,}0009 на фоне ускорения вывода в 1{,}44 раза; сравнение бэкендов выявило инверсию эффективности: WASM предпочтителен при единичных запросах ($\times$1{,}25), WebGPU при пакетной обработке.

Полученные результаты свидетельствуют о возможности эффективной реализации браузерного вывода мультиязычных трансформеров для задач анализа тональности.